\section{Functional equation, integer values, and finite products}\label{sec:functional}

The canonical product is now an independently defined object.  We derive its
recurrence rather than building the recurrence into the definition.

\subsection{Finite cancellation and the Gamma limit}

Let \(F_N(z)\) be the product in \cref{eq:Fdef} truncated after \(k=N\).
A direct cancellation gives
\begin{align}
 \frac{F_N(z)}{F_N(z-1)}
 ={}&\sqrt{2\pi}\,
 \exp\!\left(-(1+\gamma)z+\frac\gamma2-N+
                 \left(z-\frac12\right)H_N\right)\notag\\
 &\times
 \frac{(N+z)^N}{z(z+1)\cdots(z+N-1)},
 \label{eq:finite-ratio}
\end{align}
where \(H_N=\sum_{k=1}^N1/k\).  To see the only nontrivial algebra, note that
\[
 \prod_{k=1}^{N}
 \left(\frac{k+z}{k+z-1}\right)^k
 =\frac{(N+z)^N}{z\prod_{j=1}^{N-1}(z+j)}.
\]
This is a finite monomial identity and is proved by induction after clearing all
denominators.

The Gamma packet supplies the constructive Gauss limit
\begin{equation}\label{eq:gauss-gamma}
 \Gamma(z)=\lim_{N\to\infty}
 \frac{(N-1)!N^z}{z(z+1)\cdots(z+N-1)}
\end{equation}
on every compact set apart from the nonpositive integers.  Divide
\cref{eq:finite-ratio} by the finite approximation in
\cref{eq:gauss-gamma}.  The remaining factor is
\begin{equation}\label{eq:ratio-remnant}
 R_N(z)=
 \frac{\sqrt{2\pi}\,N^N}{(N-1)!N^z}
 \exp\!\left(-(1+\gamma)z+\frac\gamma2-N+
       \left(z-\frac12\right)H_N\right)
 \left(1+\frac zN\right)^N.
\end{equation}
The three standard quantitative limits are
\begin{align}
 H_N-\log N&=\gamma+O(N^{-1}),\label{eq:harmonic-limit}\\
 (N-1)!&=\sqrt{2\pi}\,N^{N-1/2}\e^{-N}(1+O(N^{-1})),
 \label{eq:stirling-fac}\\
 N\Log(1+z/N)&=z+O_R(N^{-1}).\label{eq:exp-limit}
\end{align}
Each has an explicit primitive-recursive modulus in \code{gamma\_ipra}.  On
substituting their leading terms into \cref{eq:ratio-remnant}, the powers of
\(N\) cancel, and the constant exponent is
\[
 -(1+\gamma)z+\frac\gamma2
 +(z-\tfrac12)\gamma+z=0.
\]
The quantitative error is \(O_R(1/N)\), so \(R_N(z)\to1\) uniformly on every
certified compact set.

\begin{theorem}[Functional equation]\label{thm:functional}
For every \(z\in\C\), with the identity interpreted by analytic continuation at
Gamma's poles,
\begin{equation}\label{eq:functional}
 G(z+1)=\Gamma(z)G(z),\qquad G(1)=1.
\end{equation}
On a compact zero-free domain this is a quantitative equality of fast names;
globally it is an identity of entire functions after multiplication by the
entire reciprocal-Gamma function.
\end{theorem}

\begin{proof}
Away from the nonpositive integers, combine the common product-tail modulus for
\(F_N(z)\) and \(F_N(z-1)\), the Gauss Gamma modulus, and the estimate
\(R_N(z)\to1\).  This proves \(F(z)=\Gamma(z)F(z-1)\), i.e.
\(G(z+1)=\Gamma(z)G(z)\), on a nonempty connected open domain.  Multiply by
\(1/\Gamma(z)\).  Both sides then extend to entire functions, and the equality
extends across the excluded points by their local Taylor contracts.  The
normalization is \cref{eq:G1}.
\end{proof}

For object iPRA, analytic continuation is not a black box.  At a requested
input ball intersecting a nonpositive integer, use the factorization in
\cref{thm:zeros} and the corresponding zero factor of \(1/\Gamma\).  The local
cross-multiplied identity follows from polynomial cancellation and a tail ball.
A finite cover of the input compact set combines the local proofs.

\subsection{Positive and nonpositive integer values}

Define the exact primitive-recursive sequence
\begin{equation}\label{eq:gpos}
 g_0=1,
 \qquad g_{n+1}=n!\,g_n.
\end{equation}
By \cref{eq:functional} and \(G(1)=1\),
\(G(n+1)=g_n\).  Induction yields
\begin{equation}\label{eq:integerG}
 G(n)=\prod_{j=0}^{n-2}j!
 \quad(n\ge1),
\end{equation}
where the product is empty at \(n=1\).  Thus
\[
 G(1),G(2),G(3),G(4),G(5),\ldots
 =1,1,1,2,12,288,34560,\ldots.
\]
Together with \cref{thm:zeros}, this proves
\begin{equation}\label{eq:integer-all}
 G(n)=
 \begin{cases}
 0,&n=0,-1,-2,\ldots,\\
 0!1!\cdots(n-2)!,&n=1,2,3,\ldots.
 \end{cases}
\end{equation}

\begin{theorem}[Exact iPRA integer theorem]\label{thm:integer-ipra}
Let \(\operatorname{fac}\) and \(g\) be the primitive-recursive symbols defined
by
\[
 \operatorname{fac}(0)=1,
 \quad \operatorname{fac}(n+1)=(n+1)\operatorname{fac}(n),
\]
\[
 g(0)=1,
 \quad g(n+1)=\operatorname{fac}(n)g(n).
\]
Then iPRA proves every defining equation, the superfactorial identity
\[
 g(n+1)=\prod_{j\le n}j!,
\]
and all numeral instances of the sequence above by quantifier-free induction.
\end{theorem}

\begin{proof}
The displayed recursions use the primitive-recursion constructor.  Their two
equations are instances of the zero and successor recursion axioms in
\code{PRAAx}.  Introduce a primitive-recursive accumulator for the finite
product.  Equality of its value
with \(g(n+1)\) has a quantifier-free graph formula.  The zero case is a
computation equation, and the successor case multiplies the induction
hypothesis by \((n+1)!\).  Therefore \code{PRA.qfInd} applies directly.
\end{proof}

\subsection{Superfactorials and the hyperfactorial identity}

With \(\sfac(m)=\prod_{j=0}^{m}j!\), \cref{eq:integerG} says
\begin{equation}\label{eq:superfactorial}
 G(n)=\sfac(n-2)\qquad(n\ge2),
\end{equation}
with the conventional empty extension at \(n=1\).

Let
\[
 \hyp(m)=\prod_{j=1}^{m}j^j
\]
be the hyperfactorial.  The exact finite identity on the target pages is
\begin{equation}\label{eq:hyper-cross}
 G(n)\hyp(n-1)=((n-1)!)^{n-1}\qquad(n\ge1).
\end{equation}
Indeed, count the exponent of each \(j\in\{1,\ldots,n-1\}\).  In
\(G(n)=\prod_{r=1}^{n-2}r!\), the number \(j\) appears in
\(j!, (j+1)!,\ldots,(n-2)!\), hence with exponent \(n-1-j\).  In
\(\hyp(n-1)\) it has exponent \(j\).  The total exponent is \(n-1\), which is
exactly the exponent in \(((n-1)!)^{n-1}\).

This proof can be formalized either by bounded valuation counting or, more
simply, by induction on \(n\).  The induction invariant is the division-free
identity \cref{eq:hyper-cross}; the successor step uses
\[
 G(n+1)=(n-1)!G(n),
 \quad
 \hyp(n)=\hyp(n-1)n^n,
 \quad
 n!=n(n-1)!.
\]
After substituting the induction hypothesis, both sides normalize to
\(n^n((n-1)!)^n\).  ProveIt's polynomial/power normalization modules are
designed for exactly this step.

At positive integers, \(\Gamma(n)=(n-1)!\), so
\begin{equation}\label{eq:G-over-H}
 G(n)=\frac{\Gamma(n)^{n-1}}{\hyp(n-1)}.
\end{equation}
The formal theorem is \cref{eq:hyper-cross}; the quotient form follows from the
strict positivity of \(\hyp(n-1)\).

\subsection{Finite products of Gamma values}

Iterating \cref{eq:functional} \(n\) times gives
\begin{equation}\label{eq:gamma-finite-product}
 G(k+n+1)=G(k+1)\prod_{i=1}^{n}\Gamma(k+i).
\end{equation}
Therefore, on a domain where \(G(k+1)\) is apart from zero,
\begin{equation}\label{eq:gamma-finite-quotient}
 \prod_{i=1}^{n}\Gamma(k+i)
 =\frac{G(k+n+1)}{G(k+1)}.
\end{equation}
The cross-multiplied form \cref{eq:gamma-finite-product} is global and is the
primary iPRA target.

\begin{proof}[Constructive proof]
Define
\(P_0(k)=1\) and
\(P_{n+1}(k)=P_n(k)\Gamma(k+n+1)\).  The assertion
\(G(k+n+1)=G(k+1)P_n(k)\) has the exact base case \(n=0\).  In the successor
step, multiply by \(\Gamma(k+n+1)\) and apply the functional equation once.
For named complex inputs, each multiplication is a ball operation with a
backward precision allocation.  For positive integer \(k\), every quantity is
an exact natural number, so this is ordinary quantifier-free induction.
\end{proof}

\subsection{A useful asymptotic ratio already visible from the recurrence}

For later source correction, note the exact identity for fixed \(x>0\) and
integer \(n\):
\begin{equation}\label{eq:ratio-exact}
 \frac{G(n+x)}{G(n)}
 =G(x)\,\Gamma(n)
   \prod_{j=0}^{n-1}\frac{\Gamma(x+j)}{\Gamma(j+1)}
 \qquad(n\ge1).
\end{equation}
Indeed, the extra factor \(\Gamma(n)\) exactly compensates for the last
factor introduced in the denominator product.  Gamma-ratio asymptotics, not a
power \(\Gamma(x)^n\), control the remaining product.
The precise result is proved in \cref{thm:corrected-characterization}.
